\section{The \texttt{E\_monitor} McStas Component}
Energy-sensitive monitor.

\subsection*{Identification}
\begin{itemize}
  \item \textbf{Site:} 
  \item \textbf{Author:} Kristian Nielsen and Kim Lefmann
  \item \textbf{Origin:} Risoe
  \item \textbf{Date:} April 20, 1998
\end{itemize}

\subsection*{Description}
\begin{lstlisting}
A square single monitor that measures the energy of the incoming neutrons.

Example: E_monitor(xmin=-0.1, xmax=0.1, ymin=-0.1, ymax=0.1,
Emin=1, Emax=50, nE=20, filename="Output.nrj")
\end{lstlisting}

\subsection*{Input parameters}
Parameters in \textbf{boldface} are required; the others are optional.

\begin{longtable}{p{0.22\textwidth}p{0.12\textwidth}p{0.46\textwidth}p{0.14\textwidth}}
\toprule
\textbf{Name} & \textbf{Unit} & \textbf{Description} & \textbf{Default} \\
\midrule
\endhead
nE & 1 & Number of energy channels & 20 \\
filename & string & Name of file in which to store the detector image & 0 \\
xmin & m & Lower x bound of detector opening & -0.05 \\
xmax & m & Upper x bound of detector opening & 0.05 \\
ymin & m & Lower y bound of detector opening & -0.05 \\
ymax & m & Upper y bound of detector opening & 0.05 \\
nowritefile & 1 & If set, monitor will skip writing to disk & 0 \\
xwidth & m & Width of detector. Overrides xmin, xmax. & 0 \\
yheight & m & Height of detector. Overrides ymin, ymax. & 0 \\
\textbf{Emin} & meV & Minimum energy to detect &  \\
\textbf{Emax} & meV & Maximum energy to detect &  \\
restore\_neutron & 1 & If set, the monitor does not influence the neutron state & 0 \\
\bottomrule
\end{longtable}

\subsection*{Links}
\begin{itemize}
  \item \href{run:/home/willend/willend-McCode/mcstas-comps/monitors/E_monitor.comp}{Source code} for \texttt{E\_monitor.comp}.
\end{itemize}
\IfFileExists{E_monitor_static.tex}{\input{E_monitor_static.tex}}{}